%% un-sdg.tex: UN Sustainable Development Goals (FEUP regulations)
%% -------------------------------------------------------

\chapter*{UN Sustainable Development Goals} \label{chap:UnitedNations}

The United Nations Sustainable Development Goals (SDGs) provide a global framework to achieve a better and more sustainable future for all.
It includes 17 goals to address the world's most pressing challenges, including poverty, inequality, climate change, environmental degradation, peace, and justice.

This work leverages autonomous multi-robot systems, dynamic routing algorithms, and resilient ad-hoc communication protocols (MANETs) to address these global challenges. By optimizing robotic fleet coordination and communication efficiency, this research directly contributes to technological innovation, environmental monitoring, and operational sustainability.

The specific Sustainable Development Goals mentioned have the following names:

\begin{description}
\item [SDG 7] Ensure access to affordable, reliable, sustainable and modern energy for all
\item [SDG 8] Promote sustained, inclusive and sustainable economic growth, full and productive employment and decent work for all
\item [SDG 9] Build resilient infrastructure, promote inclusive and sustainable industrialization and foster innovation
\item [SDG 11] Make cities and human settlements inclusive, safe, resilient and sustainable
\item [SDG 12] Ensure sustainable consumption and production patterns
\item [SDG 13] Take urgent action to combat climate change and its impacts
\item [SDG 15] Protect, restore and promote sustainable use of terrestrial ecosystems, sustainably manage forests, combat desertification, and halt and reverse land degradation and halt biodiversity loss
\end{description}

\begin{center}
\begin{tabular}{|l|l|p{58mm}|p{52mm}|}
\hline
\textbf{SDG} & \textbf{Target} & \textbf{Contribution} & \textbf{Performance Indicators and Metrics} \\
\hline
\hline
    \textbf{7} & 7.3 & Improving energy efficiency in robotic fleets through optimized communication scheduling (TDMA), minimizing unnecessary power consumption. & Total energy saved per mission compared to non-optimized routing protocols. \\
\hline
    \textbf{8} & 8.8 & Protecting human workers by deploying autonomous multi-robot teams to perform dangerous or repetitive tasks in hazardous environments. & Reduction in human exposure hours to high-risk or hazardous environments. \\
\hline
    \multirow{3}{*}{\textbf{9}} & 9.1 & Enhancing infrastructure flexibility by deploying mobile robots that operate via ad-hoc networks without pre-existing communication support. & Increase in efficiency and reliability of the ad-hoc network. \\
    \cline{2-4}
    & 9.5 & Leveraging advanced MANET systems and dynamic routing algorithms to solve complex coordination problems and improve operational efficiency. & Research and development metrics focused on MANET technologies. \\
    \cline{2-4}
    & 9.c & Expanding resilient communication capabilities in remote, disconnected, or industrial brownfield facilities through self-healing networks. & Network uptime and packet delivery ratio in infrastructure-less areas. \\
\hline
    \multirow{2}{*}{\textbf{11}} & 11.2 & Deploying autonomous robots for urban services such as infrastructure inspection, traffic monitoring, and emissions tracking. & Number of urban services successfully automated by autonomous robots. \\
    \cline{2-4}
    & 11.b & Enhancing community resilience to disasters by deploying rapid-response robot teams that maintain coordination when cellular networks fail. & Deployment time of the ad-hoc communication network during disruptions. \\
\hline
    \textbf{12} & 12.2 & Maximizing operational lifespan and reducing resource waste through efficient, collision-free communication scheduling that reduces energy consumption. & Decrease in energy and resource waste through optimized operations. \\
\hline
    \multirow{2}{*}{\textbf{13}} & 13.1 & Robots equipped with sensors monitor environmental conditions and provide critical data to support early warning and climate action initiatives. & Number of environmental monitoring missions completed successfully. \\
    \cline{2-4}
    & 13.3 & Assisting in disaster response, mitigation efforts, and search and rescue operations in adverse environments. & Ability of the robotic fleet to maintain group cohesion and communication. \\
\hline
    \textbf{15} & 15.1 & Deploying lightweight, coordinated swarm or mobile robots for terrestrial ecosystem monitoring, minimizing soil compaction. & Area of terrestrial ecosystem monitored per mission and reduction in ground disturbance relative to conventional survey vehicles. \\
\hline
\end{tabular}
\end{center}
